\section{One-sided residue control in ten positive variables}
\label{sec:positive-ten-centered}

We obtain a fixed polynomial in ten strictly positive integer unknowns,
of exact degree $912$ and with $126$ live binary arithmetic operations.
The inspected four-root residue predecessor has degree $1152$ and $131$
operations. Thus the degree falls by $20.83\%$, with five fewer operations.
The coefficient tuple and calibration are unchanged. This is not a
nine-variable result or a claim of global optimality.

The useful refinement is to bound a rational main-carrier residue only
\emph{from above}. Its large negative values are harmless for the positive
transport gap. An auxiliary translation and two root centers then allow a
degree-two weight instead of the predecessor's degree-$38$ majorant.
The full local proof and retained interfaces are in
\texttt{positive/ten\_centered/PROOF.md}; the argument below makes the new
inequalities and choice order explicit.

\subsection{Complete ordinary polynomial}

The fixed ports are the coherent compact44 numerals
$\beta,K_0,M_C,M_F$, with $\beta=2^{590069}-1$ and $M_F$ already shifted.
Their actual bytes and generation recipe are delivered. The only unknowns are
\[
 j,F,\alpha,i,T,X,s,n,\eta,\zeta\in\mathbb Z_{>0}.
\]
Every subsequent symbol is a polynomial abbreviation:
\begin{align*}
 J&=3j+1,\quad  q=\beta J+1,\quad  h=F+\alpha,\quad  C=q-h,\\
 Y&=q^3s,\quad  E=XY,\quad  k=\eta+\zeta,\quad  c=kY+\eta,\\
 a&=E+Y,\quad  H=4a+3,\quad  \Delta=a^2+H,\quad  \ell=EkY,\\
 R&=\bigl((q-1)(q-F)+\alpha\bigr)(q^2-1)+(M_C+qM_F)J,\\
 B_t&=(qK_0+X)C-qF,\quad  I=ic^2,\quad  D_f=1+\Delta I^2,\\
 Q&=\Delta^2I^2,\quad  \delta=Q-1,\quad  T_o=T+h,\quad  C_f=(a+2)I,\\
 V&=D_f(cT_o-R)-c,\quad D_t=\ell(\ell+k)+1,\quad D_d=1+\Delta c^2,\\
 L&=kT_o,\quad W=1+(q-1)L,\quad v=k-R-1,\\
 U&=q^5,\quad S=q^4,\quad b_0=X+ac,\quad M=(q-1)U.
\end{align*}
Set
\begin{align*}
 Z={}&M\bigl(H(\ell+k+n+L+s+V+C_f)+b_0\bigr)
       \\ &{}-HUWB_t+H(q-1)(SX-v),\\
 A&=M^2H^2D_t,\quad B=M^2H^2D_f,\\
 D&=M^2D_d,\quad C_N=M^2H^2(QV^2-1).
\end{align*}
The final division-free kernel is
\begin{align*}
 \sigma&=A+B+D,\quad \pi=AB+(A+B)D,\quad \rho=ABD,\\
 a_N&=\delta(Z^2-\sigma)+C_N,\quad t_N=\delta C_NZ^2,\\
 k_N&=a_N^2+4(t_N-\delta^2\pi),\quad r_N=16\delta^3\rho,\\
 b_N&=a_Nk_N-r_N,\\
 e_N&=k_N^2+16t_Na_N^2-4r_N(\delta Z^2+C_N),\\
 \boxed{P_{10}}&\boxed{{}=e_N^2-64t_Nb_N^2=0.}
\end{align*}
The graph uses $75$ multiplications, $34$ additions, and $17$ subtractions;
squares and scalar products are charged. Numeral leaves and reuse are free.

In proof only write $D_y=(QV^2-1)/\delta$. The exact kernel is $\delta^8$
times the product of the sixteen factors
\[
 Z-\varepsilon_1MH\sqrt{D_t}-\varepsilon_2MH\sqrt{D_f}
  -\varepsilon_3MH\sqrt{D_y}-\varepsilon_4M\sqrt{D_d}.
\]
The independent checker verifies the whole $495$-term identity.
All clearing factors are positive on the declared domain.

\subsection{Root recovery precedes the sign bootstrap}

Write $t,f,y,d$ for the positive roots of $D_t,D_f,D_y,D_d$.
The predecessor's unconditional separation proof remains valid with $T_o$
in place of $T$. Indeed $q\ge16$, $Y\ge q^3$, $c\ge2Y+1$, $a\ge2Y$, and
\[
 \ell<t<\ell+k<(a+1)c,\qquad
 f>(a+1)c^2,\qquad 0<d/H<c,\qquad t>2c+1.
\]
If the integer $cT_o-R$ is zero, then $c<y<c+1$ and the four roots are
superincreasing in order $d/H,y,t,f$. Otherwise
$y>|V|\ge D_f-c>f+t+d/H$, and their superincreasing order is $d/H,t,f,y$.
Thus every nonempty signed subset sum is nonzero. Square-class character
independence in a multiquadratic field shows that a rational signed sum
forces each root rational. At a norm zero we therefore have signed roots
$\tau_*,f_*,y_*,d_*$ and
\begin{equation}
 \ell+k+n+L+s+V+C_f-\frac{WB_t}{q-1}
 +\frac{b_0}{H}+\frac{SX-v}{U}
 =\tau_*+f_*+y_*+\frac{d_*}{H}.
 \label{eq:p10-centered-root}
\end{equation}
The first, strong and main roots are integral because their squares are
integers. The adjacent-coefficient Pell descent makes $y_*$ integral from
$QV^2-(Q-1)y_*^2=1$. It uses only integer $Q>1$ and $V$: the transformation
\[
 (V,y)\longmapsto
 ((2Q-1)V-2(Q-1)y,(2Q-1)y-2QV)
\]
preserves the reduced denominator of rational $y$, has integral first
coordinate, and strictly decreases its positive absolute value until
$|V|=1$, where $|y|=1$. No native bound or main-root orientation is used.

Integer first and main roots give Pell indices $m,p\ge1$ with
\[
 |\tau_*|=\chi_{2XY^2+1}(m),\quad k=2\psi_{2XY^2+1}(m),\qquad
 |d_*|=\chi_{a+2}(p),\quad c=\psi_{a+2}(p).
\]
Here $p\ge2$; $m=1$ would force $k=2$, $c=2Y+1<2(a+2)$, a contradiction.
Thus, \emph{before} $C>0$,
\begin{equation}
 k>8XY^2,\qquad c>8XY^3\ge8q^9.
 \label{eq:p10-centered-growth}
\end{equation}
The actual mask inequalities give the unconditional estimate
$|R|<\widehat R=q^3(q+h)$. Since $h\ge2$, \eqref{eq:p10-centered-growth}
implies $ch>\widehat R$, and therefore
\[
 0<cT_o-R<c(T+2h),\qquad V>1.
\]
As $\delta-(\Delta-1)D_f=\Delta(I^2-1)>0$ and
$\Delta-1>4(Y+1)$, the auxiliary gap satisfies
\begin{equation}
 0<y-V<\frac{V}{\delta}
 <\frac{c(T+2h)}{\Delta-1}
 <\frac{k(T+2h)}4<\frac L2.
 \label{eq:p10-centered-ybound}
\end{equation}
Moreover $f<C_f$, since $C_f^2-D_f=I^2-1>0$.

Define rational proof coordinates
\[
 r=\tau_*-\ell-k-s+\frac{WB_t}{q-1},\qquad
 \gamma=\frac{d_*-b_0}{H}-\frac{SX-v}{U}.
\]
For either main-root sign, $d_*\le d<(a+2)c$ and $2c/H<k/2$, giving
\begin{align}
 \gamma
 &<\frac{k}{2}+\frac{k}{q^5}+\frac1q+\frac{h}{q^2}
 <\frac{3k}{4}+\frac h4
 <k+\frac h2\le\frac{k(T+h)}2=\frac L2.
 \label{eq:p10-centered-gamma}
\end{align}
The middle bound uses $q\ge16$, $h\ge2$; the last uses
$(k-1)(h-1)\ge1$ and $T\ge1$. There is no lower bound on $\gamma$.
Equation~\eqref{eq:p10-centered-root} now reads
$r+\gamma+f_*+y_*=n+L+V+C_f$.
Equations~\eqref{eq:p10-centered-ybound}--\eqref{eq:p10-centered-gamma}
therefore give $r>n>0$ for all root signs. Hence
\[
 \frac{WB_t}{q-1}=\ell+k+s+r-\tau_*>0,
\]
so $B_t>0$, $C>0$, and only now $3q<R<q^4-2$ is available.

\subsection{Residues, reconstruction, and canonical completeness}

The residue argument of the four-root predecessor applies unchanged:
all added centers are integers. The normalized strong and signed auxiliary
lemmas give $p=R$. Integrality of $r+\gamma$, together with
$\gcd(H,q)=1$, yields $q^5\mid q^4X-v$. The first Pell recurrence and
$1\le m<R<q^4$ then force $R=2m-1$ and $q\mid X$; consequently
\[
 (q-1)(d_*-b_0)/H\in\mathbb Z.
\]
This step does not cancel $H$ against $q-1$. Instead Pell growth gives
$Y>q2^R$. Modulo $H$, the main Pell polynomials satisfy
\[
 \chi_{a+2}(R)-ac\equiv2^R,\qquad
 2^R(\chi_{a+2}(R)+ac)\equiv1.
\]
A negative $d_*$ would make $(q-1)(1+2^RX)/H$ an integer strictly between
zero and one. For positive $d_*$, $(q-1)(2^R-X)/H$ is an integer with
absolute value less than one. Thus $d_*>0$, $X=2^R$, and $H\mid d_*-b_0$.
Now $\gamma,r$ are integers and $(q-1)\mid B_t$.
The normalized source is restored by
\[
 w=X/q,\quad h_{\rm src}=v/(XY),\quad
 \gamma_{\rm old}=(d_*-ac-X)/H,
\]
\[
 f_{\rm old}=|f_*|,\quad \tau_{\rm old}=|\tau_*|,
 \quad y_{\rm old}=|y_*|,\quad T_{\rm src}=f_{\rm old}T_o,
 \quad z=1+\frac{(K_0+w)C-F}{q-1}.
\]
These are positive integers and satisfy every retained source equation.

Conversely choose canonical outer source witnesses and then a normalized
strong pair $f,I=ic^2$. With $M_a=\Delta I$, the canonical auxiliary
construction is
\[
 V^*=\chi_{M_a}(R)/M_a,\qquad y^*=\psi_{M_a}(R),\qquad
 T_o=\frac{V^*+c+Rf^2}{cf^2}.
\]
For $R\ge7$, its growth gives
$T_o>M_a^4/c\ge c^7>q>h$, and also $T_o>(a+2)I=C_f$.
Thus $T=T_o-h>0$. Form $L,W$ only after these choices; then $L>C_f,k$.
The canonical index gives $v\ge4XY^2$, so
\[
 \gamma=\gamma_{\rm old}-X/q+v/q^5>0.
\]
With $t_0=B_t/(q-1)\ge q$, set
$r=Wt_0+\tau-\ell-k-s>qW-2k>2L$. Finally
\[
 n=r+\gamma+f+y^*-L-V^*-C_f>r-L-C_f>0.
\]
This gives \eqref{eq:p10-centered-root} with positive roots and hence
$P_{10}=0$. No additional congruence on $Y$ or auxiliary quotient is imposed.
The choice order is outer source, strong pair, auxiliary pair, $L,W$, then
$r,n$; in particular no cofinality or freshness assumption is hidden.

\subsection{Exact degree and evidence boundary}

Let $K=MH$, and scale the ten unknowns by a formal parameter $\lambda$.
In the Laurent-series field over the rational-function coefficient field,
choose roots with the indicated leading centers. Exact defects give
\begin{align*}
 K\sqrt{D_t}&=K\ell+O(\lambda^{12}),\\
 K\sqrt{D_f}&=KC_f-\frac{KI^2}{2C_f}
                  +\text{terms of degree less than }17,\\
 K\sqrt{D_y}&=KV+O(\lambda^7),\\
 M\sqrt{D_d}&=Mb_0+O(\lambda^{11}),\\
 Z&=K(\ell+C_f+V)+Mb_0+O(\lambda^{15}).
\end{align*}
Here the large centers have degrees $21,27,49,16$, respectively; the
remaining strong-root defect has degree $17$. Of the sixteen norm factors,
eight have degree $49$, four degree $27$, two degree $21$, and two degree
$17$. Their leading terms are nonzero. Clearing $\delta^8$, of degree
$336$, gives
\[
 \deg P_{10}=8\cdot49+4\cdot27+2\cdot21+2\cdot17+336=912.
\]
Writing $q_0=3\beta j$, $k_0=\eta+\zeta$, exact simplification yields
\[
 \operatorname{in}_{912}(P_{10})=
 2^{44}q_0^{558}i^{38}k_0^{86}X^{68}s^{154}(T+F+\alpha)^8.
\]
Thus the coefficient of $j^{558}i^{38}\eta^{86}X^{68}s^{154}T^8$ is
$2^{44}(3\beta)^{558}>0$. The exact norm and circuit identities, six
actual-coefficient finite-field degree replays, and ordinary and optimized
Python runs are included. They certify formulas and measurements, not
unbounded semantics. The semantic conclusion composes the proof above
with the retained ordinary native/compiler calibration:
\[
 \mathrm{ZFC}\vdash
 \bigl(\text{$P_{10}$ has no strictly positive integer zero}\bigr)
 \longleftrightarrow\operatorname{Con}(\mathrm{ZFC}).
\]
No consistency assumption or new end-to-end formalization is claimed.
